\documentclass{lecture-kuznia}
\usepackage{textcomp}

\def\data{30.09.2025}
\def\place{Poręba Wielka}
\def\title{Funkcje tworzące}
\def\group{Finaliści}
\def\author{Autor: Stefan Świerczewski}
\def\lecturer{Prowadzący: Stefan Świerczewski}

\usepackage{xcolor}   % needed for background colors
\usepackage{mdframed} % provides the mdframed environment

\usepackage{comment}

\ifshowsolutions
  \mdfdefinestyle{solstyle}{%
    backgroundcolor=gray!10,
    linecolor=black!30,
    innertopmargin=6pt,
    innerbottommargin=6pt,
    innerleftmargin=8pt,
    innerrightmargin=8pt,
    roundcorner=6pt
  }
  \newenvironment{solutionbox}{%
    \begin{mdframed}[style=solstyle]\noindent\textbf{Rozwiązanie.}\par\smallskip
  }{%
    \end{mdframed}
  }
\else
  \excludecomment{solutionbox}
\fi


\begin{document}
\begin{center}{\Huge\textbf{\title}}\end{center}
%\LARGE{
	
	\fontsize{12}{16}\selectfont
		

        \section{\Huge Wstęp}
        Często rozwiązując zadania napotykamy ciąg postaci \(a_0, a_1, a_2 \ldots a_n, \ldots\) Nieraz trzeba wyliczyć \(k\)-ty wyraz bądź udowodnić, że wszystkie te elementy spełniają pewną zależność. Pracowanie na tych ciągach często jest żmudne, z pomocą przychodzą nam funkcje tworzące. Pozwalają nam one zapisać skomplikowane ciągi w postaci jednego obiektu na, którym możemy przeprowadzać najróżniejsze operacje.
        Zdefiniujmy zatem nasz ciąg jako współczynniki następującej funkcji(możliwie nieskończony wielomian): 
        \[f(x) = \sum_{k = 0}^{\infty}a_kx^k\]

        Zazwyczaj nie musimy się przejmować zbieżnością tego szeregu, lecz przy niektórych specyficznych problemach pozwala to ciekawe wnioski.

        W dalszej części wykładu będziemy oznaczać relację ciągu z jego funkcją tworzącą następująco:

        \[\langle a_n \rangle_{n \in \mathbb{N}} \; \sim \; A(x)\]

        Wyłuskiwanie współczynnika przy n-tej potędze będziemy oznaczać \([x^n]f(x) = a_n\)
        
        Wiemy już zatem jaki jest koncept funkcji tworzącej musimy do tego dołożyć kilka fajnych narzędzi i podstawowych funkcji.

        Jak chcemy przedstawić ciąg złożony z samych jedynek możemy powołać się na fakt, że \(1+x+x^2+\ldots+x^n = \frac{1-x^{n+1}}{1-x}\), zatem gdy \(|x| < 1\) to przechodząc do nieskończoności nasz ciąg reprezentuje funkcja:

        \[\sum_{k=0}^{\infty} x^k = \frac{1}{1-x}\]

        Otrzymaliśmy jedną z najbardziej fundamentalnych funkcji, z której często się korzysta i wyprowadza nowe funkcje tworzące.

        Częstym motywem który przetacza się w funkcjach tworzących są funkcję wymierne a w szczególności przydatna jest umiejętność rozdzielenia ich na ułamki proste.

        
        
    \section{\Huge Ważne funkcje tworzące.}

\[
1, q, q^2, q^3, \ldots = \langle q^n \rangle \sim \frac{1}{1 - qx}, \quad \text{dla } q \in \mathbb{C};
\]

\[
\binom{r}{0}, \binom{r}{1}, \binom{r}{2}, \binom{r}{3}, \ldots = \langle \binom{r}{n} \rangle \sim (1+x)^r, \quad \text{dla } r \in \mathbb{R};
\]


\[
\binom{k}{k}, \binom{k+1}{k}, \binom{k+2}{k}, \ldots = \left\langle \binom{k+n}{k} \right\rangle \sim \frac{1}{(1-x)^{k+1}}, \quad \text{dla } k \in \mathbb{N};
\]

\[
\binom{0}{k}, \binom{1}{k}, \binom{2}{k}, \ldots = \left\langle \binom{n}{k} \right\rangle \sim \frac{x^k}{(1-x)^{k+1}}, \quad \text{dla } k \in \mathbb{N};
\]

\[
1, \tfrac{1}{2}, \tfrac{1}{3}, \tfrac{1}{4}, \ldots 
= \left\langle \frac{1}{n+1} \right\rangle \sim \frac{1}{x} \ln \frac{1}{1-x};
\]

\[
1, 1, \tfrac{1}{2}, \tfrac{1}{6}, \tfrac{1}{24}, \ldots 
= \left\langle \frac{1}{n!} \right\rangle \sim e^x;
\]

\[
\langle 0, 1, 1, 2, 3, 5, \ldots \rangle = \langle F_n \rangle \sim \frac{x}{1-x-x^2}.
\]
\begin{solutionbox}
Rozważmy przykład liczb Fibonacciego i obliczmy funkcję tworzącą. Definiujemy
\[
F(x) = \sum_{n \geq 0} F_n x^n.
\]

Mamy
\[
\sum_{n \geq 1} F_{n+1} x^n
= \sum_{n \geq 1} F_n x^n + \sum_{n \geq 1} F_{n-1} x^n.
\]

Stąd
\[
\frac{F(x) - x}{x} = F(x) + xF(x),
\]
a więc
\[
F(x) = \frac{x}{1 - x - x^2}.
\]

Możemy zapisać
\[
1 - x - x^2 = (1 - r_+ x)(1 - r_- x),
\qquad r_\pm = \frac{1 \pm \sqrt{5}}{2}.
\]

Wobec tego
\[
\frac{x}{1 - x - x^2}
= \frac{1}{r_+ - r_-} \left( \frac{1}{1 - r_+ x} - \frac{1}{1 - r_- x} \right).
\]

Ponieważ
\[
\frac{1}{1 - r x} = \sum_{n \geq 0} r^n x^n,
\]
otrzymujemy
\[
F(x) = \frac{1}{\sqrt{5}} \sum_{n \geq 0} \bigl( r_+^n - r_-^n \bigr) x^n.
\]

\end{solutionbox}

    \textbf{Przydatne własności:}\\
    \begin{enumerate}
        \item Jeżeli $\langle a_n \rangle \sim A(x)$ oraz $\langle b_n \rangle \sim B(x)$, to dla dowolnych $\alpha, \beta \in \mathbb{C}$:
\[
\langle \alpha a_n + \beta b_n \rangle \sim \alpha A(x) + \beta B(x).
\]
\item Jeżeli $\langle a_n \rangle \sim A(x)$ oraz $\langle b_n \rangle \sim B(x)$, to
\[
\left\langle \sum_{k=0}^{n} a_k b_{n-k} \right\rangle \sim A(x) B(x).
\]
\item \([x^n]\{\, x^a f(x) \,\} \;=\; [x^{\,n-a}] f(x).\)
\item \([\beta x^n] f(x) \;=\; \frac{1}{\beta}\,[x^n] f(x), 
\qquad \beta \in \mathbb{R}\).
    \end{enumerate} 
    
    
  \section{\huge{Zadanka rozgrzewkowe}}

	
 
\begin{enumerate}[label=\textbf{Zad. \arabic*}]

\item Znajdź liczbę $a_n$ sposobów, w jakie można rozmieścić $n$ dolarów przy użyciu monet o nominale 1 lub 2 złotowych. 

\item Niech \(n\) i \(k\) będą liczbami całkowitymi, \(0\le k\le n\). Ile jest sposobów wyboru podzbioru dokładnie \(k\)-elementowego ze zbioru \(\{1,2,\dots,n\}\)? Wiemy, że odpowiedź to \(\binom{n}{k}\). Wyprowadźmy to korzystając z funkcji tworzących.

\begin{solutionbox}
Let \(f(n,k)\) be the number of \(k\)-element subsets of \(\{1,2,\dots,n\}\).
Partition all \(k\)-subsets into two classes:
\begin{itemize}
  \item those that contain the element \(n\), and
  \item those that do not contain \(n\).
\end{itemize}
If a \(k\)-subset contains \(n\), the remaining \(k-1\) elements must be chosen from \(\{1,\dots,n-1\}\), so there are \(f(n-1,k-1)\) such subsets. If a \(k\)-subset does not contain \(n\), then all \(k\) elements are chosen from \(\{1,\dots,n-1\}\), giving \(f(n-1,k)\) subsets. Hence the recurrence
\[
f(n,k)=f(n-1,k)+f(n-1,k-1)
\]
holds for \(0<k\le n\), with the obvious boundary conditions \(f(n,0)=1\) for all \(n\ge0\) and \(f(0,k)=0\) for \(k>0\).

Define the ordinary generating function (in the variable \(x\)) for fixed \(n\):
\[
B_n(x)=\sum_{k\ge0} f(n,k)\,x^k.
\]
Using the recurrence and summing over \(k\ge1\) we obtain
\[
\sum_{k\ge1} f(n,k)x^k
= \sum_{k\ge1} f(n-1,k)x^k + \sum_{k\ge1} f(n-1,k-1)x^k.
\]
The left-hand side equals \(B_n(x)-f(n,0)=B_n(x)-1\). The first sum on the right equals \(B_{n-1}(x)-1\), while the second sum equals \(xB_{n-1}(x)\). Therefore
\[
B_n(x)-1 = \bigl(B_{n-1}(x)-1\bigr) + xB_{n-1}(x).
\]
Rearranging yields the simple recurrence for the generating functions:
\[
B_n(x) = (1+x)\,B_{n-1}(x),\qquad B_0(x)=1.
\]
Solving this recurrence by iteration gives
\[
B_n(x)=(1+x)^n.
\]

Comparing coefficients of \(x^k\) in the binomial expansion
\[
(1+x)^n=\sum_{k=0}^n \binom{n}{k} x^k
\]
with the definition \(B_n(x)=\sum_{k\ge0} f(n,k)x^k\) we conclude
\[
f(n,k)=\binom{n}{k},
\]
as required.
\end{solutionbox}

\item Rozwiąż równania rekurencyjne
\begin{enumerate}
  \item[(a)] \[
    a_n = 2a_{n-1} - a_{n-2}, \qquad a_0 = 0,\ a_1 = 1.
  \]
  \item[(b)] \[
    b_n = 2b_{n-1} - b_{n-2} + 1, \qquad b_0 = 0,\ b_1 = 1.
  \]
  \item[(c)] \[
    c_n = 4c_{n-1} - c_{n-2}, \qquad c_0 = 1,\ c_1 = 3.
  \]
  \item[(d)] \[
    d_{n+2} = 5d_{n+1} - 6d_n - 4\cdot 2^{\,n+2}, \qquad d_0 = 1,\ d_1 = 5.
  \]
\end{enumerate}
\item Wyznacz funkcje tworzące następujących ciągów:

\begin{enumerate}
    \item $\langle 1, 0, 1, 0, 1, 0, 1, 0, \dots \rangle$;
    \item $\langle 0, -1, 2, -3, 4, -5, 6, -7, \dots \rangle$;
    \item $\langle 0, 1, 0, 3, 0, 5, 0, 7, \dots \rangle$;
    \item $\langle k, k-1, k-2, \dots, 2, 1, 0, -1, -2, \dots \rangle$, gdzie $k \in \mathbb{N}$.
\end{enumerate}
\section{\huge{Zadania}}
\item Niech $n$ będzie dodatnią liczbą całkowitą. Znajdź liczbę wielomianów $P(x)$ mających współczynniki ze zbioru \({0, 1, 2, 3}\) takich, że $P(2) = n$.
\begin{solutionbox}
Niech \(a_n\) będzie odpowiedzią na zadanie. Zdefiniujmy funkcję tworzącą
\[
f(t) = \sum_{n=0}^{\infty} a_n t^n.
\]

Niech \(P(x) = \sum_{i=0}^{k} c_i x^i\), gdzie \(c_i \in \{0,1,2,3\}\). Zauważmy, że \(P(2) = n\) wtedy i tylko wtedy, gdy
\[
\sum_{i=0}^{k} c_i 2^i = n.
\]

Ponieważ \(t^n = \prod_{i=0}^{k} t^{2^i c_i}\) i \(2^i c_i \in \{0, 2^i, 2\cdot 2^i, 3\cdot 2^i\}\), funkcję \(f(t)\) można rozłożyć w postaci
\[
f(t) = \prod_{i=0}^{\infty} \left(1 + t^{2^i} + t^{2\cdot 2^i} + t^{3\cdot 2^i}\right)
= \prod_{i=0}^{\infty} \frac{1 - t^{4 \cdot 2^i}}{1 - t^{2^i}}
= \frac{1}{1-t} \cdot \frac{1}{1-t^2}.
\]

Rozwijając w szereg otrzymujemy
\[
\frac{1}{(1-t)(1-t^2)} = \frac{1}{2} \left( \frac{1}{(1-t)^2} + \frac{1}{1-t^2} \right)
= \frac{1}{2} \sum_{k=0}^{\infty} k t^{k-1} + \sum_{k=0}^{\infty} t^{2k}
= \sum_{n=0}^{\infty} \left\lfloor \frac{n}{2} \right\rfloor + 1 \, t^n.
\]

Stąd liczba takich wielomianów wynosi
\[
a_n = \left\lfloor \frac{n}{2} \right\rfloor + 1.
\]

\end{solutionbox}

\item Niech \(p\) będzie nieparzystą liczbą pierwszą. Ile jest podzbiorów \(A\) zbioru 
\(\{1, 2, \dots, 2p\}\) o \(p\) elementach, dla których suma elementów jest podzielna przez \(p\)?
\begin{solutionbox}

Niech $\omega$ będzie pierwotną $p$-tą pierwiastkiem jednostki. Wówczas
\[
\prod_{k=1}^{2p} (x - \omega^k) = (x^p - 1)^2 = x^{2p} - 2x^p + 1.
\]

Porównując współczynnik przy $x^p$, mamy
\[
2 = \sum_{\{k_1+k_2+\dots+k_p\}} \omega^{k_1 + k_2 + \dots + k_p} = \sum_{j=0}^{p-1} n_j \omega^j,
\]
gdzie pierwsza suma jest po wszystkich podzbiorach $\{k_1, k_2, \dots, k_p\}$ zbioru $\{1,2,\dots,2p\}$, a $n_j$ w drugiej sumie jest liczbą takich podzbiorów, dla których
\[
k_1 + k_2 + \dots + k_p \equiv j \pmod{p}.
\]

Zatem $\omega$ jest pierwiastkiem wielomianu
\[
G(x) = (n_0 - 2) + \sum_{j=1}^{p-1} n_j x^j,
\]
który jest wielomianem stopnia $p-1$. Przypomnijmy, że minimalny wielomian $\omega$ nad $\mathbb{Q}[x]$ jest
\[
F(x) = \sum_{j=0}^{p-1} x^j,
\]
również stopnia $p-1$. Ponieważ minimalne wielomiany są jednoznaczne do stałego współczynnika, $G(x)$ jest stałą wielokrotnością $F(x)$. To implikuje, że
\[
n_0 - 2 = n_1 = n_2 = \dots = n_{p-1}.
\]

Skoro $\sum_{j=0}^{p-1} n_j = \binom{2p}{p}$, to
\[
n_0 = \frac{1}{p} \binom{2p}{p} - 2 + 2 = \frac{1}{p} \binom{2p}{p}.
\]
\end{solutionbox}

\section{\huge{Szczypta teorii}}

Niech $\langle a_n \rangle_{n \in \mathbb{N}} \overset{e}{\sim} A_e(x)$ oznacza, że $A_e(x)$ jest wykładniczą funkcją tworzącą ciągu $\langle a_n \rangle$, czyli szeregiem formalnym określonym następująco:
\[
A_e(x) = \sum_{n \in \mathbb{N}} a_n \frac{x^n}{n!}.
\]

\textbf{(Splot dwumianowy)} Jeżeli 
\[
\langle a_n \rangle \overset{e}{\sim} A_e(x) \quad \text{oraz} \quad \langle b_n \rangle \overset{e}{\sim} B_e(x),
\]
to
\[
\sum_{k=0}^{n} \binom{n}{k} a_k b_{n-k} \overset{e}{\sim} A_e(x) B_e(x).
\]

\section{\huge{Zadania bardzo trudne}}
\item Niech $n$ i $k$ będą liczbami całkowitymi dodatnimi, przy czym $k \ge n$ oraz $k - n$ jest liczbą parzystą. Rozważamy $2n$ lamp oznaczonych $1,2,\dots,2n$, z których każda może być w stanie włączonym lub wyłączonym. Początkowo wszystkie lampy są wyłączone. Rozważamy sekwencję kroków, w której w każdym kroku zmieniana jest jedna z lamp (z włączonego na wyłączoną lub z wyłączonej na włączoną).

Niech $N$ oznacza liczbę takich sekwencji składających się z $k$ kroków, które prowadzą do stanu, w którym lampy $1$ do $n$ są wszystkie włączone, a lampy $n+1$ do $2n$ są wszystkie wyłączone.

Niech $M$ oznacza liczbę takich sekwencji składających się z $k$ kroków, które prowadzą do stanu, w którym lampy $1$ do $n$ są wszystkie włączone, lampy $n+1$ do $2n$ są wszystkie wyłączone, a żadna z lamp $n+1$ do $2n$ nie była nigdy włączona.

Wyznacz stosunek $\frac{N}{M}$.
\begin{solutionbox}
Niech
\[
	f(x)=\left(x+\frac{x^3}{3!}+\frac{x^5}{5!}+\cdots\right)
	\left(1+\frac{x^2}{2!}+\frac{x^4}{4!}+\cdots\right)
	=\frac{\sinh(2x)}{2}.
\]
Wtedy $N=k!\,[x^k](f(x)^n)$. Z kolei dla $g(x)=\sinh(x)$ mamy
$M=k!\,[x^k](g(x)^n)$. Stąd
\[
	\frac{N}{M}
	=2^{-n}\frac{[x^k](\sinh(2x))^n}{[x^k](\sinh(x))^n}
	=2^{-n}\cdot 2^k
	=2^{k-n}.
\]
\end{solutionbox}

\end{enumerate}
\end{document}
